\section{Which Distinctions Matter?}
\label{sec:main-value}

Adding a distinction has a canonical minimal form. For $D=U^{-1}(1)$,
\[
\Pi\vee U
=
\{B\cap D,B\cap D^c:B\in\Pi\}\setminus\{\varnothing\}.
\]

\begin{proposition}[Minimal join]
\label{prop:main-minimal-join}
$\Pi\vee U$ is the unique coarsest refinement of $\Pi$ that makes $U$ measurable.
\end{proposition}

The operation splits exactly the old atoms crossed by the new distinction and leaves all other atoms unchanged. To see minimality, let $\calQ$ be any refinement of $\Pi$ that makes $D=U^{-1}(1)$ measurable. Every atom $Q\in\calQ$ lies inside some old atom $B$ and, because $D$ is $\calQ$-measurable, lies either in $D$ or in $D^c$. Hence $Q$ lies inside either $B\cap D$ or $B\cap D^c$, so $\calQ$ refines $\Pi\vee U$. This matters operationally: refinement should not be equated with wholesale model growth. The canonical information update adds the least observability required by the proposed bit. Its value is task-dependent:
\[
V_L(U\mid\Pi)=R_L(\Pi)-R_L(\Pi\vee U).
\]
Risk monotonicity gives $V_L(U\mid\Pi)\ge0$ at the population optimum, because every old predictor remains feasible after refinement. Finite implementations can still overfit, which is why Section~\ref{sec:main-guard} uses held-out scoring and penalties before committing a split.

\begin{theorem}[Loss-relevant refinement]
\label{thm:main-loss-relevant}
Let $\Pi'$ refine $\Pi$. Assume Bayes actions exist and the $\Pi'$-Bayes action is unique on positive-mass atoms. Then
\[
R_L(\Pi)=R_L(\Pi')
\]
if and only if the $\Pi'$-Bayes action has a $\Pi$-measurable version. Equivalently, in the finite partition model, zero value means the refined Bayes action is constant across all refined atoms inside each old atom $B\in\Pi$.
\end{theorem}

\begin{proof}[Proof sketch]
If the refined Bayes action is $\Pi$-measurable, it is feasible under $\Pi$, so $R_L(\Pi)\le R_L(\Pi')$; monotonicity gives equality. If risks are equal, any $\Pi$-Bayes action is also $\Pi'$-optimal. Uniqueness under $\Pi'$ forces agreement with the refined Bayes action on positive-mass atoms.
\end{proof}

For finite $Y$ under log loss, $R_{\log}(\Pi)=H(Y\mid\Pi)$. The Bayes action in atom $B$ is the empirical or population conditional law $P(Y\mid B)$, and the risk is the conditional entropy averaged over atoms. For a binary distinction $U$,
\[
V_{\log}(U\mid\Pi)
=
R_{\log}(\Pi)-R_{\log}(\Pi\vee U)
=
I(Y;U\mid\Pi).
\]
Equivalently,
\[
V_{\log}(U\mid\Pi)
=
\sum_{B\in\Pi}P(B)
\sum_{u\in\{0,1\}}P(U=u\mid B)
D_{\mathrm{KL}}\!\left(
P(Y\mid B,U=u)\,\Vert\,P(Y\mid B)
\right).
\]
Thus all candidate bits may have the same marginal frequency while only consequence-relevant bits have positive conditional mutual information. This identity is the reason BalancedMask-HQS is a useful diagnostic: the stream makes all query bits equally frequent, so frequency or geometric novelty cannot identify the correct bit, while $I(Y;U\mid\Pi)$ can.

For $Y\in L^2(\Omega;\R^d)$ and squared loss,
\[
R_2(\Pi)-R_2(\Pi\vee U)
=
\E\left[
\left\|
\E[Y\mid \Pi,U]
-
\E[Y\mid \Pi]
\right\|_2^2
\right].
\]
We use standard value-of-information identities here, specialized to the refine-or-repair online setting \citep{cover2006elements,gneiting2007strictly}. The value of a distinction is not geometric novelty; it is consequence value. Consequently, a visually clean cluster split can be worthless if both children have the same Bayes action, while a low-level bit can be valuable if it changes the conditional consequence. This is also why the algorithm compares refinement against repair rather than accepting every statistically nonzero split: a split can have positive value but still be smaller than the gain from simply repairing stale payloads.
